\documentclass[12pt]{article}
\usepackage{amsmath,amssymb}
\usepackage{geometry}
\usepackage{enumitem}
\geometry{margin=0.8in}
\setlength{\parindent}{0pt}

\begin{document}

\begin{center}
  {\Large\textbf{Algebra I: Exponent Rules}}\\[3pt]
  {\large Timed Practice --- 30 Minutes}
\end{center}

Name: \rule{2.7in}{0.4pt}\hfill Date: \rule{1.5in}{0.4pt}

\section*{Instructions}
Complete all 18 problems in 30 minutes without a calculator. You may use this
reference page. Show key steps and give exact answers. Write simplified
answers with positive exponents only; do not leave fractions within fractions.
Assume all variables in expressions being simplified are positive real numbers.
In Problem 16, solve for real $x$.

\begingroup
\small
\section*{Fundamental Exponent Rules}
For the integer-exponent rules below, $m,n\in\mathbb Z$ and all expressions
must be defined. In particular, denominators and bases raised to zero or
negative powers must be nonzero.

\medskip
\begingroup
\renewcommand{\arraystretch}{1.7}
\begin{tabular}{@{}p{1.9in}p{4.7in}@{}}
  \textbf{Rule} & \textbf{Formula}\\
  First power & $a^1=a$\\
  Product of powers & $a^m\cdot a^n=a^{m+n}$\\
  Quotient of powers & $\displaystyle\frac{a^m}{a^n}=a^{m-n}$\\
  Power of a power & $(a^m)^n=a^{mn}$\\
  Power of a product & $(ab)^n=a^nb^n$\\
  Power of a quotient & $\displaystyle\left(\frac{a}{b}\right)^n=\frac{a^n}{b^n}$\\
  Zero exponent & $a^0=1\quad(a\ne0)$\\
  Negative exponent & $\displaystyle a^{-n}=\frac{1}{a^n},\qquad
    \left(\frac{a}{b}\right)^{-n}=\left(\frac{b}{a}\right)^n$\\
  Rational exponent & $\displaystyle a^{m/n}=\sqrt[n]{a^m}
    =\left(\sqrt[n]{a}\right)^m$\\
\end{tabular}
\endgroup

For the rational-exponent rule, take $a>0$, $m\in\mathbb Z$, and positive
integer $n$. The product, quotient, and power rules also hold for rational
exponents when all bases are positive.

\section*{Important Conditions and Conventions}
\begin{itemize}[leftmargin=*,itemsep=0.35em]
  \item For positive integers $n$, $0^n=0$. The expressions $0^0$ and $0^{-n}$
  are undefined here.
  \item A negative exponent means a reciprocal, not a negative value.
  \item For positive integers $n$, $(-a)^n=a^n$ when $n$ is even and
  $(-a)^n=-a^n$ when $n$ is odd. In $-a^n$, the exponent applies to $a$ only.
  \item Exponents distribute over products and quotients, not sums or
  differences. In general, $(a+b)^n\ne a^n+b^n$.
\end{itemize}
\endgroup

\newpage
\begin{center}
  {\large\textbf{Exponent Rules Practice}}\\[2pt]
  Problems 1--6: approximately 6 minutes
\end{center}

Simplify each expression.
\begin{enumerate}[leftmargin=*,itemsep=0pt]
  \item $x^7\cdot x^4$
  \vspace{\stretch{1}}

  \item $\displaystyle\frac{y^{12}}{y^5}$
  \vspace{\stretch{1}}

  \item $(a^3)^4$
  \vspace{\stretch{1}}

  \item $(-3x^2y)^3$
  \vspace{\stretch{1}}

  \item $\displaystyle\left(\frac{4a^3}{3b^2}\right)^2$
  \vspace{\stretch{1}}

  \item $(7x^3y^2)^0+2^0$
  \vspace*{\stretch{1}}
\end{enumerate}

\newpage
\begin{center}
  {\large\textbf{Exponent Rules Practice --- Continued}}\\[2pt]
  Problems 7--12: approximately 10 minutes
\end{center}

Simplify each expression.
\begin{enumerate}[leftmargin=*,itemsep=0pt,resume]
  \item $\displaystyle\left(\frac{5}{2}\right)^{-2}$
  \vspace{\stretch{1}}

  \item $-2^4+(-2)^4+(-2)^3$
  \vspace{\stretch{1}}

  \item $\displaystyle\frac{12x^{-3}y^5}{3x^2y^{-1}}$
  \vspace{\stretch{1}}

  \item $(2a^{-2}b^3)^{-2}$
  \vspace{\stretch{1}}

  \item $\displaystyle
    \frac{(-3m^2n^{-1})^2(2m^{-3}n^4)}{6m^{-1}n^2}$
  \vspace{\stretch{1}}

  \item $\displaystyle\left(\frac{4x^3y^{-2}}{2x^{-1}y}\right)^{-2}$
  \vspace*{\stretch{1}}
\end{enumerate}

\newpage
\begin{center}
  {\large\textbf{Exponent Rules Practice --- Continued}}\\[2pt]
  Problems 13--18: approximately 14 minutes
\end{center}

\begin{enumerate}[leftmargin=*,itemsep=0pt,resume]
  \item Evaluate exactly: $\displaystyle16^{3/4}+27^{2/3}$.
  \vspace{\stretch{1}}

  \item Write $x^{5/2}$ in simplified radical form, and write
  $\sqrt[3]{a^4}$ using a rational exponent.
  \vspace{\stretch{1}}

  \item Simplify:
  $\displaystyle\frac{(x^{3/2}y^{-1/3})^2}{x^{-1}y^{1/3}}$.
  \vspace{\stretch{1}}

  \item Solve: $\displaystyle3^{2x-1}=\frac{27^x}{9}$.
  \vspace{\stretch{1}}

  \item A student writes
  $\displaystyle(2a^2b^{-1})^3=\frac{6a^5}{b^3}$.
  Explain each error and give the correct simplification.
  \vspace{\stretch{1}}

  \item Simplify $\displaystyle\frac{x^{-2}+y^{-2}}{x^{-1}y^{-1}}$,
  then evaluate your result at $x=2$, $y=3$.
  \vspace*{\stretch{1}}
\end{enumerate}

\end{document}
